\bigrafa{revisar negritas en las defs}

In this section, we define the standard notions of complexity theory used throughout this thesis.

\begin{definition}[Language]
    In this work we consider the following notion of problems: given a sequence of input tokens, output a token as the answer. 
    Mathematically, given a vocabulary $\Sigma$, we call a mapping $\gL:\cup_{k\in\mathbb{N}^+}\Sigma^k\to\Sigma$ a problem. 
    If the correct answer is always $\qaccept$ or $\qreject$, we call $\gL$ a decision problem.
    In circuit complexity, such $\gL$ is also called a \textbf{language}.
\end{definition}

\subsection{Boolean circuits}

The Boolean circuits are a model of computation generalizing Boolean formulas.
It is a simplified model of the silicon chips used to make modern computers. 
It is a natural model for nonuniform computation, as is the case for transformers; thus, most of our results on the expressiveness of transformers are stated in the language of circuit theory.


A Boolean circuit is a procedural diagram showing how to derive an output from a binary input string by applying a sequence of basic Boolean operations $\OR$ ($\lor$), $\AND$ ($\land$), and $\NOT$ ($\lnot$) on the input bits.

Though the standard definition of circuit complexity only deals with binary strings, given any finite vocabulary $\Sigma$, we can always replace each token in $\Sigma$ by its binary representation, and the length of the input only blows up by a constant factor. 
Therefore, we can extend existing complexity classes to arbitrary finite vocabulary naturally. 

\begin{definition}[Boolean circuits] 
    For every $n \in \mathbb{N}$, an $n$-input, single-output Boolean circuit is a directed acyclic graph with $n$ sources (vertices with no incoming edges) and one sink (vertex with no outgoing edges). 
    All nonsource vertices are called gates and are labeled with one of $\land$, $\lor$ or $\lnot$ (i.e., the logical operations $\OR$, $\AND$, and $\NOT$).
    The vertices labeled with $\lor$ and $\land$ have fan-in (i.e., number of incoming edges) equal to $2$ and the vertices labeled with $\lnot$ have fan-in $1$. 
    The size of $C$, denoted by $|C|$, is the number of vertices in it.

    If $C$ is a Boolean circuit, and $x \in \{0,1\}^n$ is some input, then the output of $C$ on $x$, denoted by $C(x)$, is defined in the natural way. 
    More formally, for every vertex $v$ of $C$, we give it a value $val(v)$ as follows: 
    If $v$ is the $i$-th input vertex then $val(v) = x_i$ and otherwise $val(v)$ is defined recursively by applying $v$'s logical operation on the values of the vertices connected to $v$. 
    The output $C(x)$ is the value of the output vertex.
\end{definition}


\begin{definition}[Circuit families and language recognition] 
    Let $T : \mathbb{N} \to \mathbb{N}$ be a function. 
    A $T(n)$-size circuit family is a sequence $\{C_n\}_{n\in\mathbb{N}}$ of Boolean circuits, where $C_n$ has $n$ inputs and a single output, and its size $|C_n| \leq T(n)$ for every $n$.
    We say that a language $\gL$ is in $\SIZE(T(n))$ if there exists a $T(n)$-size circuit family $\{C_n\}_{n\in\mathbb{N}}$ such that for every $x\in\{0,1\}^n$, $x \in \gL \iff C_n(x)=1$.
\end{definition}

Interesting complexity classes arise when we consider “small” circuits, such as the class of languages recognized by all polynomial-size families.

\begin{definition}[$\ppoly$]
    $\ppoly$ is the class of languages that are decidable by polynomial-sized circuit families.
    That is, $\ppoly = \cup_c \SIZE(n^c)$.
\end{definition}

The class $\ppoly$ fits rather awkwardly in the complexity world since it contains undecidable languages such as the language Unary halt\footnote{UHALT = \{$1^n$ : $n$'s binary expansion encodes a pair $\langle M, x \rangle$ such that $M$ halts on input $x$\}}. 
The root of the problem is that for a language $\gL$ to be in $\ppoly$, it suffices that a circuit family for $\gL$ exists even if we have no way of actually constructing the circuits.
Thus, it may be fruitful to try to restrict attention to circuits that can actually be built, say using a fairly efficient Turing machine.

\begin{definition}[$\Ptime$-uniform circuit families]
    A circuit family $\{C_n\}_{n\in\mathbb{N}}$ is $\Ptime$-uniform if there is a polynomial-time Turing machine that on input $1^n$ outputs the description of the circuit $C_n$.
\end{definition}

\bigskip

Boolean circuits are a good model to represent parallel computation.
The \textbf{depth} of a circuit is the length of the longest directed path from an input node to the output node.

\begin{definition}[$\NC$]
    For every $d$, a language $\gL$ is in $\NC^d$ if $\gL$ can be decided by a family of circuits $\{C_n\}_{n\in\mathbb{N}}$ where $C_n$ has $\poly(n)$ size and depth $O(\log^d(n))$. 
    The class $\NC$ is $\cup_{i \geq 1}\NC^i$.
\end{definition}

When studying the depth of a circuit, the fan-in of the gates becomes of interest.
A polynomial fan-in can be simulated with a logarithmic blow-up in depth.
This motivates the study of circuits with unbounded fan-in.

\begin{definition}[$\AC$]
    The class $\AC^i$ is defined similarly to $\NC^i$ except gates are allowed to have unbounded fan-in (i.e., the $\OR$ and $\AND$ gates can be applied to more than two bits). 
    The class $\AC$ is $\cup_{i \geq 0}\AC^i$.
\end{definition}

$\AC^0$ circuits allow for an interesting number of operations.
For example, addition and multiplication of numbers of fixed precision can be done with $\AC^0$ circuits.
Also, the equality of two strings of bits can be done in $\AC^0$.
From the ability to compute equality it emerges that any function $f$ with $O(\log(n))$ bits as input can be computed in $\AC^0$.
The circuit computing $f$ enumerates the result for each of the inputs and outputs the one that correspond to the branch equal to the input.
Let $k$ be the number of bits of the input of $f$, then $f$ can be represented as the following formula: $\land_{i = 0}^{2^k} f(i) \land i = x$.
Observe that the formula represents a circuit of constant depth.
Also note that the formula has size $O(2^k)$; then, if $k \in O(\log(n))$ the circuit has polynomial size; and, therefore, it is in $\AC^0$.


One of the main results of circuit complexity is that the $\MAJORITY$\footnote{$\MAJORITY = \left\{x : \frac{\sum_{i=1}^{|x|} x_i}{|x|} > \frac{1}{2}\right\}$, i.e., the words with more $1$s than $0$s} problem cannot be computed by $\AC^0$ circuits.
\rafa{cita a algún lugar donde prueben eso o es lo suficientemente sabido?}

\begin{theorem}\label{maj-notin-ac0}
    $\MAJORITY \notin \AC^0$
\end{theorem}

This theorem motivates the study of threshold circuits, circuits allowed to have $\MAJORITY$ gates.

\begin{definition}[$\TC$]
    The class $\TC^i$ is defined similarly to $\AC^i$ except the circuits can have gates $\OR$, $\AND$, $\NOT$ and $\MAJORITY$ of unbounded fan-in.
    The class $\TC$ is $\cup_{i \geq 0}\TC^i$.
\end{definition}


Following Theorem \ref{maj-notin-ac0}, a family of variations of $\MAJORITY$ began to be studied.

\begin{definition}[$\MAJORITY$ with promise]
    For every $e \in (1/2, 1]$ we define the problem \textbf{Majority with promise} ($\PROMISEDMAJORITY{e}$) as $\MAJORITY$ but restricted to words where there are at least $e$ $0$s or $e$ 1s.
\end{definition}

An interesting result from descriptive complexity is that for all $e$, there is a circuit in $\AC^0$ recognizing $\PROMISEDMAJORITY{e}$\rafa{agregar cita}.

\subsubsection{Randomized Boolean circuits}

In this thesis we also work with randomized Boolean circuits.
These are Boolean circuits, but they are allowed to have random gates.
These gates have fan-in $0$, and their value is a uniformly distributed bit.
This value is fixed once the evaluation of the circuit starts.
Then, the evaluation of a randomized circuit $C$ with input $x$ can be thought as evaluating a Boolean circuit $C'$ with input $x r$, where $r$ are the values of the random gates.

\begin{definition}[Randomized Boolean circuits]
    A \textbf{randomized Boolean circuit} is a directed acyclic graph with $n+r$ sources and one sink.
    All nonsource vertices are labeled as $\land$, $\lor$ or $\lnot$ (i.e., the logical operations $\OR$, $\AND$, and $\NOT$).
    The vertices labeled with $\lor$ and $\land$ have fan-in (i.e., number of incoming edges) equal to $2$ and the vertices labeled with $\lnot$ have fan-in $1$. 
    The size of $C$, denoted by $|C|$, is the number of vertices in it.
    The source vertices are either labeled as $\mathsf{input}$ or $\mathsf{random}$.
    
    If $C$ is a Boolean circuit, and $x \in \{0,1\}^n$ is some input, then the output of $C$ on $x$, denoted by $C(x)$, is defined as random variable in the following way.
    For every vertex $v$ of $C$, we give it a value $val(v)$ as follows: 
    If $v$ is the $i$-th input vertex, then $val(v) = x_i$; if $v$ is labeled as $\mathsf{random}$, then $val(v)$ is a uniform random Boolean variable; and otherwise $val(v)$ is defined recursively by applying $v$'s logical operation on the values of the vertices connected to $v$. 
    The output $C(x)$ is the value of the output vertex.
\end{definition}

In general, we say that a family of random circuits recognizes a language $\gL$ when $\Pr[C_{|x|}(x) = \gL(x)]$ is high.

\begin{definition}[Randomized circuit families and language recognition] 
    Let $T : \mathbb{N} \to \mathbb{N}$ be a function.
    A $T(n)$-size randomized circuit family is a sequence $\{C_n\}_{n\in\mathbb{N}}$ of randomized Boolean circuits, where $C_n$ has $n$ inputs and a single output, and its size $|C_n| \leq T(n)$ for every $n$.
    We say that a language $\gL$ is in $\BPSIZE(T(n))[f(n)]$ if there exists a $T(n)$-size randomized circuit family $\{C_n\}_{n\in\mathbb{N}}$ with $O(f(n))$ random gates such that for every $x\in\{0,1\}^n$, $\Pr[C_n(x) = 1] \geq 2/3 \iff \gL(x)$.
\end{definition}

Since every random Boolean circuit can be interpreted as Boolean circuit with the random bits padded in the input, we can also define recognition in the following way:

\begin{definition}
    We say that a language $\gL$ is in $\BPSIZE(T(n))[f(n)]$ if there exists a $T(n)$-size circuit family $\{C'_n\}_{n\in\mathbb{N}}$ such that for every $x\in\{0,1\}^n$, $\Pr_{r \in \{0,1\}^{f(n)}}[C'_n(xr) = 1] \geq 2/3 \iff \gL(x)$.
\end{definition}

In general, we drop the bound for the number of random bits when it is not of interest.

\begin{definition}
    $\BPSIZE(T(n)) = \BPSIZE(T(n))[T(n)]$.
\end{definition}

\begin{definition}
    Similar to $\ppoly$, we define $\BPppoly$ as the class of languages that are decidable by polynomial-sized random circuit families.
    That is, $\BPppoly = \cup_c \BPSIZE(n^c)$.
\end{definition}

We extend the definition of $\NC$, $\NC^i$, $\AC$, $\AC^i$, $\TC$, $\TC^i$ uniform and not uniform to randomized Boolean circuits in the natural way.
We denote by $\mathcal{C}[f(n)]$ the class $\mathcal{C}$ restricted to $O(f(n))$ random gates.

Observe that a Boolean circuit is a random Boolean circuit without random gates.

\begin{theorem}
    Every circuit class $\mathcal{C}$ is contained in its corresponding randomized circuit class.
    For example $\AC^0 \subseteq \BPAC^0$.
\end{theorem}

The other inclusion is not known in general.
In the case of $\AC^0$ and $\BPAC^0$, restricting $\BPAC^0$ to only $O(f(n))$ random bits is enough to be make it fall inside $\AC^0$.
This is because in $\AC^0$ we can enumerate all possible values for the random bits.
Then, we compute the output of the random Boolean circuit for each case.
There are $O(n)$ possible random values.
We can run the circuit for each of them in parallel leading only to a linear blow-up in width but constant in depth.
Then, by definition of the random circuit, at least $2/3$ of the cases have the same output.
Therefore, we can output $\PROMISEDMAJORITY{2/3}$.

\begin{theorem}\label{const-random-in-ac0}
    $\BPAC^0[f(n)] \subseteq \AC^0$
\end{theorem}


\subsection{Boolean circuits and Turing machines}

Most of the work in complexity theory is done over Turing machines.
In this section, we present some known relations between circuits and Turing machines.

\begin{definition}[$\DTIME$]
    Let $T : \mathbb{N} \to \mathbb{N}$ be some function. 
    A language $\gL$ is in $\DTIME(T(n))$ iff there is a Turing machine that runs in time $c \cdot T(n)$ for some constant $c > 0$ and decides $\gL$.
\end{definition}

\begin{definition}
    $\Ptime = \cup_{c \geq 0} \DTIME(n^c)$.
\end{definition}

A probabilistic Turing machine is a Turing machine with access to a random coin.

\begin{definition}[$\BPTIME$]
    Let $T : \mathbb{N} \to \mathbb{N}$ be some function. 
    A language $\gL$ is in $\BPTIME(T(n))$ iff there is a probabilistic Turing machine $M$ that runs in time $c \cdot T(n)$ for some constant $c > 0$ and for all $x \in \{0,1\}^n$, $x \in \gL \iff \Pr[M(x) = 1] \geq 2/3$
\end{definition}

\begin{definition}
    $\BPP = \cup_{c \geq 0} \BPTIME(n^c)$.
\end{definition}

When restricting $\ppoly$ and $\BPppoly$ to only their uniform families they coincide with $\Ptime$ and $\BPP$.
To simulate a circuit by a Turing machine, one first computes the circuit description (this can be done since we are only analyzing uniform families) and then evaluates it, both in polynomial time.
For the simulation of Turing machines by circuits, the construction of Cook-Levin works.

\begin{theorem} $\mathsf{Uniform} \cap \ppoly = \Ptime$ \end{theorem}

\begin{theorem} $\mathsf{Uniform} \cap \BPppoly = \BPP$ \end{theorem}

An important result relating machines with and without randomness is Adleman's theorem.
Adleman proved that the non-uniformity of $\ppoly$ gives enough power to simulate the randomness of $\BPP$. \rafa{agregar cita a Adleman}

\begin{theorem}[Adleman's theorem]\label{adleman}
    $\BPP \subseteq \ppoly$
\end{theorem}



{\color{orange}
Lista de cosas que voy referenciando:

\begin{itemize}
    \item $\MAJORITY_pr$ con $pr > 1/2$ se puede hacer en ac0 (https://www.ccs.neu.edu/home/viola/papers/BPvsE.pdf)
\end{itemize}
}